%% ============================================================================
%% Template LaTeX baseado no artigo publicado em:
%% Elsevier - Results in Engineering (Review Article, 2 colunas)
%% Classe: elsarticle
%% Compilacao recomendada: pdflatex / xelatex / tectonic
%% ============================================================================
%% DICA DE FORMATO:
%% - Para visual identico ao PDF publicado (duas colunas):
\documentclass[final,5p,times,twocolumn]{elsarticle}
%% - Para submissao em formato preprint (uma coluna, 12pt), comente a linha
%%   acima e descomente a linha abaixo:
%% \documentclass[preprint,12pt]{elsarticle}

%% --- Pacotes essenciais ---
\usepackage[utf8]{inputenc}
\usepackage[T1]{fontenc}
\usepackage{amsmath,amssymb,amsfonts}
\usepackage{graphicx}
\usepackage{booktabs}
\usepackage{tabularx}
\usepackage{multirow}
\usepackage{array}
\usepackage{microtype}
\usepackage{lineno}
\usepackage[colorlinks=true,linkcolor=blue,citecolor=blue,urlcolor=blue]{hyperref}

%% --- Configuracao de citacoes numericas (padrao do artigo: [1], [6-22]) ---
\biboptions{sort&compress}

%% --- Configuracao da revista ---
\journal{Results in Engineering}

%% --- Ajuste opcional de numeracao de linhas (util no modo preprint) ---
\modulolinenumbers[5]

%% --- Configuracao de colunas flexiveis em tabelas ---
\newcolumntype{L}[1]{>{\raggedright\arraybackslash}p{#1}}
\newcolumntype{Y}{>{\raggedright\arraybackslash}X}

\begin{document}

%% ============================================================================
%% CABECALHO (FRONTMATTER)
%% ============================================================================
\begin{frontmatter}

%% Tipo de artigo (opcional como nota de titulo, como "Review article")
\title{[Article Title: A Comprehensive Review of Techniques and Innovations]}

%% --- Autores e correspondencia (padrao a, b, c... igual ao artigo) ---
\author[a,b]{[First Author]\corref{cor1}}
\ead{first.author@institution.edu}

\author[c]{[Second Author]\corref{cor1}}
\ead{second.author@institution.edu}

\author[a]{[Third Author]}

\author[d]{[Fourth Author]}

\author[e]{[Fifth Author]\corref{cor1}}
\ead{fifth.author@institution.edu}

\author[f]{[Sixth Author]}

\cortext[cor1]{Corresponding authors.}

%% --- Afiliacoes ---
\affiliation[a]{%
    organization={[School/Department Name, Faculty Name, University Name]},
    city={[City]},
    postcode={[Postcode]},
    country={[Country]}
}

\affiliation[b]{%
    organization={[Department Name, Faculty Name, University Name]},
    city={[City]},
    country={[Country]}
}

\affiliation[c]{%
    organization={[Department Name, University Name]},
    postcode={[Postcode]},
    country={[Country]}
}

\affiliation[d]{%
    organization={[School/Department Name, University Name]},
    postcode={[Postcode]},
    country={[Country]}
}

\affiliation[e]{%
    organization={[School/Department Name, University Name]},
    city={[City]},
    postcode={[Postcode]},
    country={[Country]}
}

\affiliation[f]{%
    organization={[School/Department Name, University Name]},
    postcode={[Postcode]},
    country={[Country]}
}

%% --- Abstract ---
\begin{abstract}
[Single-paragraph abstract of approximately 200--250 words. Summarize the scope of the review, the systematic taxonomy established to evaluate the methodologies, the main comparative findings and trade-offs, recent mathematical or methodological advancements, emerging research directions, unresolved challenges, and concluding trajectories for future research.]
\end{abstract}

%% --- Keywords ---
\begin{keyword}
[Keyword 1] \sep
[Keyword 2] \sep
[Keyword 3] \sep
[Keyword 4] \sep
[Keyword 5]
\end{keyword}

\end{frontmatter}

%% Descomente a linha abaixo se desejar numeracao de linhas na revisao:
%% \linenumbers

%% ============================================================================
%% SUMARIO (CONTENTS) — Caracteristico de artigos de revisao na revista
%% ============================================================================
\setcounter{tocdepth}{2} % Mostra secoes e subsecoes (1, 2, 2.1, 2.2...)
\tableofcontents
\vspace{1em}

%% ============================================================================
%% SECAO 1: INTRODUCTION
%% ============================================================================
\section{Introduction}\label{sec:intro}

[Opening paragraph introducing the engineering problem, real-world relevance, and motivation, with numeric citations \cite{ref1,ref2}. Refer to a full-width introductory figure as shown in Fig.~\ref{fig:overview}.]

[Second paragraph explaining why direct measurement or conventional solutions are impractical, and defining the fundamental problem categories (e.g., forward problem vs.\ inverse problem), shown schematically in Fig.~\ref{fig:schematic}.]

[Third paragraph categorizing the main classes of approaches in the literature \cite{ref1,ref2,ref3}, contrasting their main strengths and drawbacks, see Fig.~\ref{fig:comparison}.]

[Fourth paragraph summarizing secondary approaches and referencing a comparative summary table, see Table~\ref{tab:summary_methods}.]

%% Exemplo de figura em largura total (duas colunas) — equivalente a Fig. 1
\begin{figure*}[t]
    \centering
    %% Substitua o \fbox abaixo por: \includegraphics[width=0.9\textwidth]{fig1.pdf}
    \fbox{\parbox[c][4.5cm][c]{0.88\textwidth}{\centering [Figure 1: Full-width illustrative figure spanning both columns]}}
    \caption{[Descriptive caption for Fig.~1, explaining the main components or risk zones highlighted \cite{ref3}.]}
    \label{fig:overview}
\end{figure*}

%% Exemplo de tabela em largura total (duas colunas) — equivalente a Table 1
\begin{table*}[t]
    \centering
    \caption{[Illustrative summary of the main approaches reviewed in Section~\ref{sec:intro}.]}
    \label{tab:summary_methods}
    \small
    \begin{tabularx}{\textwidth}{L{2.8cm} Y Y Y}
        \toprule
        \textbf{[Approach Category]} & \textbf{[Example]} & \textbf{[Shortcoming]} & \textbf{[Solution]} \\
        \midrule
        \textit{Category 1} &
        (i) [Method A], \newline (ii) [Method B], \newline (iii) [Method C] &
        (i) [Limitation 1], \newline (ii) [Limitation 2], \newline (iii) [Limitation 3] &
        [Proposed solution using, e.g., (i) Technique A, (ii) Technique B] \\
        \addlinespace
        \textit{Category 2} &
        [Method D] &
        [Limitation of Method D in complex or anisotropic scenarios] &
        [Reformulate as a nonlinear problem and solve via optimization] \\
        \addlinespace
        \textit{Category 3} &
        (i) [Database approach], \newline (ii) [Neural Networks] &
        [Need large training dataset] &
        [Artificially generate or augment dataset] \\
        \bottomrule
    \end{tabularx}
\end{table*}

%% Exemplo de figura em coluna unica — equivalente a Fig. 2
\begin{figure}[t]
    \centering
    %% Substitua por: \includegraphics[width=\columnwidth]{fig2.pdf}
    \fbox{\parbox[c][3.2cm][c]{0.92\columnwidth}{\centering [Figure 2: Single-column schematic diagram]}}
    \caption{[Caption for single-column schematic diagram (e.g., Forward vs.\ inverse problem).]}
    \label{fig:schematic}
\end{figure}

%% Exemplo de figura em coluna unica — equivalente a Fig. 3
\begin{figure}[t]
    \centering
    %% Substitua por: \includegraphics[width=\columnwidth]{fig3.pdf}
    \fbox{\parbox[c][3.8cm][c]{0.92\columnwidth}{\centering [Figure 3: Classification and comparison diagram]}}
    \caption{[Classification and comparison of different algorithms reviewed in the paper.]}
    \label{fig:comparison}
\end{figure}

[Paragraph defining the theoretical background (e.g., well-posedness vs.\ ill-posedness conditions) and the main contributions of this review paper.]

The paper is organized as follows. Section~\ref{sec:sec2} presents [topic of Section 2]. Next, [topic of Section 3] is presented in Section~\ref{sec:sec3}. A review of [topic of Section 4] is presented in Section~\ref{sec:sec4}. Section~\ref{sec:sec5} introduces [topic of Section 5]. Section~\ref{sec:sec6} discusses [topic of Section 6]. Finally, concluding remarks are provided in Section~\ref{sec:conclusions}.

%% ============================================================================
%% SECAO 2: MODEL-BASED APPROACH (PART I)
%% ============================================================================
\section{[Model-based approach I: Reconstruction]}\label{sec:sec2}

[Overview of the first major family of model-based techniques, introducing the main sub-approaches and referencing Fig.~\ref{fig:sec2_pros}.]

\begin{figure*}[t]
    \centering
    %% Substitua por: \includegraphics[width=0.85\textwidth]{fig4.pdf}
    \fbox{\parbox[c][3.2cm][c]{0.85\textwidth}{\centering [Figure 4: Taxonomy / pros of each model-based approach]}}
    \caption{[The pros of each model-based approach for the target problem.]}
    \label{fig:sec2_pros}
\end{figure*}

\subsection{[Primary technique: Deconvolution]}\label{subsec:sec2_1}

[Introduction to the first sub-approach, its main limitations, and its subdivision into three domains, see Fig.~\ref{fig:sec2_class}.]

\begin{figure}[t]
    \centering
    %% Substitua por: \includegraphics[width=0.75\columnwidth]{fig5.pdf}
    \fbox{\parbox[c][3.5cm][c]{0.75\columnwidth}{\centering [Figure 5: Classification of sub-approaches]}}
    \caption{[Classification of approaches within Section~\ref{subsec:sec2_1}.]}
    \label{fig:sec2_class}
\end{figure}

\subsubsection{[Domain A: Time domain formulation]}\label{subsubsec:sec2_1_1}

[Formulation of the continuous model, illustrated in Fig.~\ref{fig:sec2_workflow}. The relationship between response $r$ and input $f$ is expressed as:]
\begin{equation}\label{eq:convolution}
    r(y,t) = \int_{0}^{t} T_{s}(x,y,t-\zeta)\, f(x,\zeta)\, \mathrm{d}\zeta,
\end{equation}
where $T_{s}(x,y,t-\zeta)$, $s = 1, \dots, n$, denotes the transfer function. Discretizing Eq.~\eqref{eq:convolution} yields:
\begin{equation}\label{eq:discrete}
    r = T_{s} f,
\end{equation}
where $r \in \mathbb{R}^{m}$, $f \in \mathbb{R}^{m}$, and $T_{s} \in \mathbb{R}^{m \times m}$ are defined as:
\begin{equation}\label{eq:matrices}
    r =
    \begin{bmatrix}
        r(\Delta t) \\
        r(2\Delta t) \\
        \vdots \\
        r(m\Delta t)
    \end{bmatrix},
    \quad
    f =
    \begin{bmatrix}
        f(\Delta t) \\
        f(2\Delta t) \\
        \vdots \\
        f(m\Delta t)
    \end{bmatrix}.
\end{equation}
The solution of Eq.~\eqref{eq:discrete} can be formulated as the least-squares problem:
\begin{equation}\label{eq:lsq}
    \min \| r - T_{s} f \|_{2}^{2}.
\end{equation}

\begin{figure*}[t]
    \centering
    %% Substitua por: \includegraphics[width=0.9\textwidth]{fig6.pdf}
    \fbox{\parbox[c][3.0cm][c]{0.88\textwidth}{\centering [Figure 6: Workflow illustration spanning two columns]}}
    \caption{[Illustration of the reconstruction workflow (inputs, transfer functions, and output).]}
    \label{fig:sec2_workflow}
\end{figure*}

\subsubsection{[Domain B: Frequency domain formulation]}\label{subsubsec:sec2_1_2}

[Spectral formulation obtained via Fourier Transform of Eq.~\eqref{eq:convolution}:]
\begin{equation}\label{eq:freq}
    \tilde{r}(y,\omega) = \tilde{T}_{s}(x,y,\omega)\,\tilde{f}(x,\omega),
\end{equation}
with the pseudo-inverse solution given by:
\begin{equation}\label{eq:pinv}
    \tilde{f}(\omega) = \tilde{T}_{s}(\omega)^{+} \tilde{r}(\omega),
    \qquad
    \tilde{T}_{s}^{+} = (\tilde{T}_{s}^{H}\tilde{T}_{s})^{-1}\tilde{T}_{s}^{H}.
\end{equation}

\subsubsection{[Domain C: Wavelet / Transform formulation]}\label{subsubsec:sec2_1_3}

[Description of finite-duration basis expansion. Example of paired subequations `(7a)` and `(7b)` matching the article style:]
\begin{subequations}\label{eq:wavelet_basis}
\begin{align}
    \psi_{m,n}(\tau) &= a_{0}^{-m/2}\psi\!\left(\frac{\tau - n b_{0} a_{0}^{m}}{a_{0}^{m}}\right), \label{eq:wavelet_a} \\
    \varphi_{M,n}(\tau) &= a_{0}^{-M/2}\varphi\!\left(\frac{\tau - n b_{0} a_{0}^{M}}{a_{0}^{M}}\right). \label{eq:wavelet_b}
\end{align}
\end{subequations}

\subsection{[Secondary technique: State estimation-based approach]}\label{subsec:sec2_2}

[Introduction to state-space formulation, starting from the linear equation of motion:]
\begin{equation}\label{eq:eom}
    M\ddot{u}(t) + C\dot{u}(t) + K u(t) = f(t),
\end{equation}
and its continuous-time state-space representation:
\begin{equation}\label{eq:ss}
    \dot{x}(t) = A x(t) + B f(t), \qquad r(t) = G x(t) + J f(t).
\end{equation}

\subsubsection{[Sub-method A: Markov parameters computation]}\label{subsubsec:sec2_2_1}

[Discussion and discrete convolution representation using forward and inverse Markov parameters.]

\subsubsection{[Sub-method B: Kalman filter approaches]}\label{subsubsec:sec2_2_2}

[Overview of deterministic-stochastic methods, followed by structured item descriptions:]
\begin{itemize}
    \item \textbf{[Variant 1 (e.g., Augmented Kalman filter)]:} [Brief description of formulation, strengths, and limitations.]
    \item \textbf{[Variant 2 (e.g., Dual Kalman filter)]:} [Brief description of formulation, strengths, and limitations.]
    \item \textbf{[Variant 3 (e.g., Sparse Kalman filter)]:} [Brief description of formulation, strengths, and limitations.]
    \item \textbf{[Variant 4 (e.g., Unscented / Extended Kalman filter)]:} [Brief description of formulation, strengths, and limitations.]
\end{itemize}

[Short wrap-up paragraph comparing the model-based methods reviewed in Section~\ref{sec:sec2}.]

%% ============================================================================
%% SECAO 3: MODEL-BASED APPROACH (PART II)
%% ============================================================================
\section{[Model-based approach II: Localization]}\label{sec:sec3}

[Brief introduction linking Section~\ref{sec:sec2} and Section~\ref{sec:sec3}.]

\subsection{[Methods based on time difference and geometry]}\label{subsec:sec3_1}

[Explanation of wave-propagation and geometry-based strategies (e.g., triangulation), illustrated in Fig.~\ref{fig:triangulation}.]

\begin{figure}[t]
    \centering
    %% Substitua por: \includegraphics[width=0.9\columnwidth]{fig7.pdf}
    \fbox{\parbox[c][3.5cm][c]{0.88\columnwidth}{\centering [Figure 7: Geometric localization schematic]}}
    \caption{[Geometric method for localizing the target source.]}
    \label{fig:triangulation}
\end{figure}

\subsection{[Methods based on structural vibrations]}\label{subsec:sec3_2}

[Overview of vibration-based methods, divided into similarity searching and optimization-based techniques.]

\subsubsection{[Similarity searching methods]}\label{subsubsec:sec3_2_1}

[Introductory paragraph followed by bulleted sub-topics with equations and figures (Fig.~\ref{fig:crosscorr} and Fig.~\ref{fig:cosine}):]

\begin{itemize}
    \item \textbf{[Method 1 (e.g., Time reversal method)]:} [Formulation and center-of-gravity localization equations.]
    \item \textbf{[Method 2 (e.g., Normalized cross-correlation)]:} [Cross-correlation definition and reference database workflow, see Fig.~\ref{fig:crosscorr}.]
    \item \textbf{[Method 3 (e.g., Cosine similarity technique)]:} [Cosine metric definition and piecewise reference signal:]
    \begin{equation}\label{eq:piecewise}
        p(t) =
        \begin{cases}
            \sin\!\left(\dfrac{\pi}{a} t\right), & 0 < t < a, \\
            0, & a < t < T.
        \end{cases}
    \end{equation}
    \item \textbf{[Method 4 (e.g., Wavelet indexing strategy)]:} [Continuous wavelet transform similarity index.]
\end{itemize}

\begin{figure*}[t]
    \centering
    %% Substitua por: \includegraphics[width=0.9\textwidth]{fig8.pdf}
    \fbox{\parbox[c][4.2cm][c]{0.88\textwidth}{\centering [Figure 8: Full-width reference database and signal comparison diagram]}}
    \caption{[General concept of the similarity-based localization method using a reference database.]}
    \label{fig:crosscorr}
\end{figure*}

\begin{figure}[t]
    \centering
    %% Substitua por: \includegraphics[width=0.75\columnwidth]{fig9.pdf}
    \fbox{\parbox[c][2.8cm][c]{0.75\columnwidth}{\centering [Figure 9: Vector similarity illustration]}}
    \caption{[Geometric interpretation of the similarity metric between two vectors.]}
    \label{fig:cosine}
\end{figure}

\subsubsection{[Optimization-based methods]}\label{subsubsec:sec3_2_2}

[Objective function formulation, e.g., $E = \sum_{i=1}^{m} \| T_{s}^{(i)} f - r^{(i)} \|$, and discussion of gradient-based vs.\ metaheuristic algorithms (GA, PSO).]

%% ============================================================================
%% SECAO 4: MACHINE AND DEEP LEARNING APPROACH
%% ============================================================================
\section{[Machine and deep learning approach]}\label{sec:sec4}

[Review of ML and DL architectures applied to the problem, highlighting their advantages, limitations, and physics-informed hybrid frameworks.]

\begin{enumerate}
    \item \textbf{[Key Point / Strategy 1]:} [Explanation of input representation or feature engineering.]
    \item \textbf{[Key Point / Strategy 2]:} [Explanation of physics-based data augmentation or architecture design.]
    \item \textbf{[Key Point / Strategy 3]:} [Explanation of attention mechanisms, incremental learning, or explainability.]
\end{enumerate}

[Discussion referencing the literature summary in Table~\ref{tab:ml_papers}.]

%% Exemplo de tabela ampla de revisao da literatura — equivalente a Table 2
\begin{table*}[t]
    \centering
    \caption{[A summary of representative papers using machine learning and deep learning models.]}
    \label{tab:ml_papers}
    \small
    \begin{tabularx}{\textwidth}{L{1.2cm} Y Y Y}
        \toprule
        \textbf{Ref.} & \textbf{Main aim} & \textbf{Technical model} & \textbf{Conclusion} \\
        \midrule
        \cite{ref1} & [Objective of study 1] & [Model 1 (e.g., SVM)] & [Main finding or comparative accuracy] \\
        \addlinespace
        \cite{ref2} & [Objective of study 2] & [Model 2 (e.g., GNN / CNN)] & [Main finding or comparative accuracy] \\
        \addlinespace
        \cite{ref3} & [Objective of study 3] & [Model 3 (e.g., ANN + GA)] & [Main finding or comparative accuracy] \\
        \bottomrule
    \end{tabularx}
\end{table*}

\subsection{[Machine learning data requirements]}\label{subsec:sec4_1}

[Discussion of data scarcity and data-efficient learning paradigms:]
\begin{enumerate}
    \item \textbf{Synthetic data generation:} [Role of physics-based simulations such as FEM.]
    \item \textbf{Data augmentation:} [Signal transformations to improve generalization.]
    \item \textbf{Transfer learning:} [Domain adaptation from simulation to experimental targets.]
    \item \textbf{Self-supervised and unsupervised learning:} [Anomaly detection without baseline labels.]
    \item \textbf{Zero-shot and few-shot learning:} [Embedding physical laws for rare events.]
\end{enumerate}

%% ============================================================================
%% SECAO 5: MATHEMATICAL TECHNIQUES TO DEAL WITH ILL-POSEDNESS
%% ============================================================================
\section{[Mathematical techniques to deal with ill-posedness]}\label{sec:sec5}

[Unified notation $r = H f$ and overview of stabilization strategies against measurement noise.]

\subsection{[Posedness investigation]}\label{subsec:sec5_1}

[Singular Value Decomposition (SVD) and condition number analysis:]
\begin{equation}\label{eq:svd}
    H = U \Sigma V^{T}, \qquad \operatorname{cond}(H) = \frac{\sigma_{1}}{\sigma_{m}}.
\end{equation}

\subsection{[The least squares solution]}\label{subsec:sec5_2}

[Objective function $J = \| H f - r \|_{2}$, normal equations, and sensitivity to small singular values $\sigma_{i}$.]

\subsection{[Sensor placement and measurement optimization]}\label{subsec:sec5_3}

[Discussion of sensor count, coherence-based placement strategies, and reduction of the condition number.]

\subsection{[Regularization methods]}\label{subsec:sec5_4}

[Classification into direct and iterative regularization approaches, summarized in Table~\ref{tab:regularization}.]

\subsubsection{[Direct approach for regularization]}\label{subsubsec:sec5_4_1}

[Review of Tikhonov regularization (L-curve, GCV), TSVD, Bayesian regularization, sparse ($l_{1}$) regularization, and hybrid methods.]

\subsubsection{[Iterative regularization]}\label{subsubsec:sec5_4_2}

[Review of Levenberg--Marquardt, Landweber, Kaczmarz, and Krylov subspace methods (CG, LSQR, LSMR), followed by a summary of common weaknesses:]
\begin{enumerate}[(i)]
    \item [Convergence to local minima / suboptimal performance.]
    \item [Computational cost for large-scale iterations.]
    \item [Sensitivity to algorithm-specific parameters.]
    \item [Slow convergence in strongly ill-conditioned problems.]
    \item [Dependence on the quality of the initial guess.]
\end{enumerate}

%% Exemplo de tabela comparativa agrupada com multirow — equivalente a Table 3
\begin{table*}[t]
    \centering
    \caption{[A comparative summary of direct and iterative regularization methods.]}
    \label{tab:regularization}
    \small
    \begin{tabularx}{\textwidth}{L{1.8cm} L{3.0cm} Y Y}
        \toprule
        \multicolumn{2}{l}{\textbf{Regularization method}} & \textbf{Strength} & \textbf{Weakness} \\
        \midrule
        \multirow{4}{1.8cm}{\textbf{Direct approach}}
            & Tikhonov regularization
            & [Simplicity and global smoothness]
            & (i) [Ineffective for large-scale problems] \newline (ii) [Not suitable for non-smooth signals] \\
        \cmidrule(lr){2-4}
            & Truncated SVD (TSVD)
            & [Effective against high-frequency instability]
            & [Lower identification accuracy in general] \\
        \cmidrule(lr){2-4}
            & Bayesian regularization
            & (i) [Adaptive parameter selection] \newline (ii) [Uncertainty quantification]
            & [Possible non-physical fluctuations] \\
        \cmidrule(lr){2-4}
            & Sparse regularization
            & [Effective for sparse signals in time/space]
            & (i) [Noise sensitivity] \newline (ii) [Computational complexity] \\
        \midrule
        \multirow{4}{1.8cm}{\textbf{Iterative approach}}
            & Levenberg--Marquardt
            & \multirow{4}{=}{(i) [Suitable for large-scale inverse problems] \newline (ii) [Higher accuracy and robustness without explicit SVD]}
            & \multirow{4}{=}{(i) [Suboptimal convergence risk] \newline (ii) [Computational cost] \newline (iii) [Parameter sensitivity] \newline (iv) [Slow convergence] \newline (v) [Initial guess dependence]} \\
            & Landweber & & \\
            & Kaczmarz & & \\
            & Krylov subspace & & \\
        \bottomrule
    \end{tabularx}
\end{table*}

%% ============================================================================
%% SECAO 6: APPLICABILITY AND REAL-TIME FEASIBILITY
%% ============================================================================
\section{[Applicability and real-time feasibility]}\label{sec:sec6}

[Critical discussion of latency, computational load, hardware constraints, hybrid physics--data approaches, and practical deployment in real-time monitoring and active control systems.]

%% ============================================================================
%% SECAO 7: CONCLUSIONS
%% ============================================================================

%% ============================================================================
%% SECAO EXPERIMENTAL: DATASETS, METODOS E RESULTADOS DE SEGMENTACAO
%% Inserida para o estudo de segmentacao de campos fisicos 2D do The Well.
%% As celulas com \textemdash{} devem ser preenchidas apos a execucao.
%% ============================================================================

\section{Experimental protocol and comparative results}
\label{sec:experimental}

\subsection{Datasets and physical segmentation targets}
\label{subsec:datasets}

Table~\ref{tab:datasets} summarizes the ten two-dimensional datasets selected
from The Well. The physical target and the corresponding ground-truth rule
must be specified for each dataset before training. The proposed targets below
are candidates and must be validated against the variables available in each
dataset; they are not pre-existing annotation labels.

\begin{table*}[t]
    \centering
    \caption{Selected The Well datasets and candidate segmentation targets.}
    \label{tab:datasets}
    \small
    \begin{tabularx}{\textwidth}{L{3.15cm} L{1.65cm} L{2.35cm} Y Y}
        \toprule
        \textbf{Dataset} & \textbf{Grid} & \textbf{Physical system} &
        \textbf{Candidate target} & \textbf{Ground-truth definition / variable} \\
        \midrule
        \texttt{acoustic\_scattering} & $256\times256$ & Acoustic scattering &
        Acoustic-field regions & [Specify field, physical criterion and threshold] \\
        \addlinespace
        \texttt{active\_matter} & $256\times256$ & Active matter &
        Active or organized structures & [Specify order/activity variable and criterion] \\
        \addlinespace
        \texttt{euler\_multi\_quadrants} & $512\times512$ & Compressible flow &
        Density/flow structures & [Specify physical variable(s) and criterion] \\
        \addlinespace
        \texttt{gray\_scott\_reaction\_diffusion} & $128\times128$ & Reaction--diffusion &
        Concentration patterns & [Specify species/concentration and criterion] \\
        \addlinespace
        \texttt{helmholtz\_staircase} & $1024\times256$ & Wave propagation &
        Wave-field regions & [Specify field amplitude/energy and criterion] \\
        \addlinespace
        \texttt{planetswe} & $256\times512$ & Shallow-water dynamics &
        Flow structures & [Specify coordinate-aware physical variable and criterion] \\
        \addlinespace
        \texttt{rayleigh\_benard} & $512\times128$ & Thermal convection &
        Convective structures & [Specify temperature/velocity-derived criterion] \\
        \addlinespace
        \texttt{shear\_flow} & $128\times256$ & Shear flow &
        Vortical/shear regions & [Specify vorticity or shear criterion] \\
        \addlinespace
        \texttt{turbulent\_radiative\_layer\_2D} & $128\times384$ & Radiative turbulent flow &
        Thermodynamic/flow regions & [Specify physical variable(s) and criterion] \\
        \addlinespace
        \texttt{viscoelastic\_instability\_v2} & $512\times512$ & Viscoelastic flow &
        Deformation/flow structures & [Specify stress, strain or flow criterion] \\
        \bottomrule
    \end{tabularx}
\end{table*}

\subsection{Pre-processing methods}
\label{subsec:preprocessing}

The pre-processing methods evaluated in this study are listed in
Table~\ref{tab:preprocessing}. Parameters must be estimated using the training
partition only and then kept fixed for validation and test data. The unfiltered
reference condition should use a documented, consistent input scaling.

\begin{table*}[t]
    \centering
    \caption{Pre-processing methods and parameters to be reported.}
    \label{tab:preprocessing}
    \small
    \begin{tabularx}{\textwidth}{L{0.65cm} L{3.0cm} L{2.35cm} Y Y}
        \toprule
        \textbf{ID} & \textbf{Method} & \textbf{Category} &
        \textbf{Purpose} & \textbf{Parameters to report} \\
        \midrule
        P1 & Z-score normalization & Numerical scaling &
        Standardize each input channel & [Training mean and standard deviation] \\
        P2 & Min--Max normalization & Numerical scaling &
        Map values to a defined interval & [Minimum, maximum and target interval] \\
        P3 & Robust normalization & Numerical scaling &
        Reduce sensitivity to extreme values & [Median and interquartile range] \\
        P4 & Percentile clipping & Numerical scaling &
        Limit extreme values & [Lower and upper percentiles] \\
        P5 & Logarithmic transformation & Numerical transform &
        Reduce positive-value distribution skewness & [Transform, offset and variable(s)] \\
        P6 & Gaussian filter & Spatial filtering &
        Suppress high-frequency variation & [Standard deviation $\sigma$ and kernel size] \\
        P7 & Median filter & Spatial filtering &
        Reduce impulsive noise & [Window size] \\
        P8 & Bilateral filter & Spatial filtering &
        Smooth while retaining selected edges & [Spatial and intensity parameters] \\
        P9 & Anisotropic diffusion & Spatial filtering &
        Smooth while preserving selected structures & [Iterations, conductance and step size] \\
        P10 & Wavelet denoising & Multiscale filtering &
        Suppress selected wavelet coefficients & [Wavelet family, level and threshold rule] \\
        \bottomrule
    \end{tabularx}
\end{table*}

\subsection{Segmentation architectures}
\label{subsec:architectures}

Table~\ref{tab:architectures} summarizes the three segmentation architectures.
For a fair comparison, the input channels, output classes, loss function,
training schedule, data split and model-selection rule must be documented.
For multi-channel physical fields, the input stem of each model must be adapted
to the number of physical variables used.

\begin{table*}[t]
    \centering
    \caption{Deep-learning architectures included in the comparison.}
    \label{tab:architectures}
    \small
    \begin{tabularx}{\textwidth}{L{2.4cm} L{2.55cm} Y Y}
        \toprule
        \textbf{Architecture} & \textbf{Family} & \textbf{Role in the study} &
        \textbf{Implementation details to report} \\
        \midrule
        U-Net & Convolutional encoder--decoder &
        Convolutional baseline for spatial segmentation &
        [Depth, channels, kernel sizes, normalization and parameter count] \\
        \addlinespace
        SegFormer & Hierarchical Transformer &
        Transformer-based semantic segmentation comparison &
        [Variant/backbone, input adaptation, decoder and parameter count] \\
        \addlinespace
        Mask2Former & Mask-classification Transformer &
        Mask-based segmentation comparison &
        [Backbone, number of queries, decoder settings and parameter count] \\
        \bottomrule
    \end{tabularx}
\end{table*}

\subsection{Post-processing methods}
\label{subsec:postprocessing}

Table~\ref{tab:postprocessing} lists the post-processing methods. Each method
must operate on a clearly specified representation (probability map, binary
mask, multiclass label map or temporal sequence). Thresholds and structural
parameters must be selected using validation data only.

\begin{table*}[t]
    \centering
    \caption{Post-processing methods and parameters to be reported.}
    \label{tab:postprocessing}
    \small
    \begin{tabularx}{\textwidth}{L{0.65cm} L{3.0cm} L{2.5cm} Y Y}
        \toprule
        \textbf{ID} & \textbf{Method} & \textbf{Input representation} &
        \textbf{Purpose} & \textbf{Parameters to report} \\
        \midrule
        Q1 & Probability thresholding & Probability map &
        Convert probabilities into class labels & [Threshold per class] \\
        Q2 & Morphological opening & Binary/class mask &
        Remove selected small foreground structures & [Structuring element and size] \\
        Q3 & Morphological closing & Binary/class mask &
        Close selected small gaps & [Structuring element and size] \\
        Q4 & Opening followed by closing & Binary/class mask &
        Combine removal of protrusions and gap filling & [Operation order and element size] \\
        Q5 & Small-component removal & Connected-component map &
        Remove components below an area criterion & [Connectivity and minimum area] \\
        Q6 & Hole filling & Binary/class mask &
        Fill enclosed background regions & [Connectivity and hole-size limit] \\
        Q7 & Connected-component selection & Component labels &
        Retain components according to a predefined rule & [Connectivity and selection rule] \\
        Q8 & Watershed & Probability/distance or marker map &
        Separate touching regions where physically justified & [Markers, connectivity and compactness] \\
        Q9 & Conditional Random Field (CRF) & Unary scores and pairwise features &
        Refine spatial label consistency & [Unary/pairwise weights and iterations] \\
        Q10 & Temporal filtering & Probability/mask sequence &
        Evaluate temporal consistency across frames & [Temporal window and filter rule] \\
        \bottomrule
    \end{tabularx}
\end{table*}

\subsection{Experimental design and result reporting}
\label{subsec:experimental_design}

The full factorial design contains ten pre-processing conditions, three
architectures and ten post-processing conditions for each dataset. For ten
datasets, this corresponds to $10\times3\times10\times10=3000$ evaluated
combinations. Since post-processing is applied to saved model predictions,
it does not require retraining the segmentation network for every
post-processing method. The principal training count is therefore
$10\times3\times10=300$ runs, excluding reference baselines, repetitions and
hyperparameter-search runs.

Table~\ref{tab:full_results} is the long-format results template. Each row
represents one dataset--pre-processing--architecture--post-processing
combination. Duplicate the row for each evaluated combination, or store the
complete results in a supplementary CSV file and report selected summaries
in the article.

\begin{table*}[t]
    \centering
    \caption{Full experimental results: one row per evaluated combination.}
    \label{tab:full_results}
    \scriptsize
    \setlength{\tabcolsep}{3pt}
    \begin{tabularx}{\textwidth}{L{2.05cm} L{0.55cm} L{1.45cm} L{1.5cm} c c c c c}
        \toprule
        \textbf{Dataset} & \textbf{ID} & \textbf{Pre-processing} &
        \textbf{Architecture / post-processing} &
        \textbf{IoU} & \textbf{Dice} & \textbf{Prec.} & \textbf{Rec.} &
        \textbf{Latency (ms)} \\
        \midrule
        \textit{Dataset} & [P/Q] & [Method] & [Model / method] &
        [\,] & [\,] & [\,] & [\,] & [\,] \\
        \bottomrule
    \end{tabularx}
    \vspace{2pt}
    \begin{minipage}{\textwidth}
        \scriptsize
        \textit{Note:} Report mean and standard deviation across independent
        runs or test trajectories. Specify whether latency includes
        post-processing and report the hardware used.
    \end{minipage}
\end{table*}

Table~\ref{tab:architecture_comparison} compares the architectures while
aggregating over the selected pre- and post-processing conditions. The
aggregation rule must be stated (e.g., mean over all valid combinations or
performance under the validation-selected configuration).

\begin{table*}[t]
    \centering
    \caption{Architecture comparison by dataset.}
    \label{tab:architecture_comparison}
    \small
    \begin{tabularx}{\textwidth}{L{3.0cm} *{3}{>{\centering\arraybackslash}X}}
        \toprule
        \textbf{Dataset} & \textbf{U-Net} & \textbf{SegFormer} & \textbf{Mask2Former} \\
        \midrule
        \texttt{acoustic\_scattering} & [IoU / Dice] & [IoU / Dice] & [IoU / Dice] \\
        \texttt{active\_matter} & [IoU / Dice] & [IoU / Dice] & [IoU / Dice] \\
        \texttt{euler\_multi\_quadrants} & [IoU / Dice] & [IoU / Dice] & [IoU / Dice] \\
        \texttt{gray\_scott\_reaction\_diffusion} & [IoU / Dice] & [IoU / Dice] & [IoU / Dice] \\
        \texttt{helmholtz\_staircase} & [IoU / Dice] & [IoU / Dice] & [IoU / Dice] \\
        \texttt{planetswe} & [IoU / Dice] & [IoU / Dice] & [IoU / Dice] \\
        \texttt{rayleigh\_benard} & [IoU / Dice] & [IoU / Dice] & [IoU / Dice] \\
        \texttt{shear\_flow} & [IoU / Dice] & [IoU / Dice] & [IoU / Dice] \\
        \texttt{turbulent\_radiative\_layer\_2D} & [IoU / Dice] & [IoU / Dice] & [IoU / Dice] \\
        \texttt{viscoelastic\_instability\_v2} & [IoU / Dice] & [IoU / Dice] & [IoU / Dice] \\
        \bottomrule
    \end{tabularx}
\end{table*}

Table~\ref{tab:prepost_effect} isolates the contribution of the
pre-processing and post-processing stages. Report the change relative to a
clearly defined baseline, evaluated on the same test cases.

\begin{table*}[t]
    \centering
    \caption{Ablation analysis of pre-processing and post-processing effects.}
    \label{tab:prepost_effect}
    \small
    \begin{tabularx}{\textwidth}{L{2.4cm} L{2.3cm} L{2.5cm} Y c c}
        \toprule
        \textbf{Dataset} & \textbf{Architecture} & \textbf{Stage / method} &
        \textbf{Baseline} & \textbf{IoU change} & \textbf{Dice change} \\
        \midrule
        \textit{Dataset} & [Architecture] & [P1--P10 or Q1--Q10] &
        [Reference condition] & [\,] & [\,] \\
        \bottomrule
    \end{tabularx}
\end{table*}

Table~\ref{tab:best_configuration} reports the best configuration selected
using validation data and evaluated once on the held-out test set. The
selection criterion (e.g., validation Dice) must be fixed before test
evaluation.

\begin{table*}[t]
    \centering
    \caption{Best validation-selected configuration and held-out test performance.}
    \label{tab:best_configuration}
    \small
    \begin{tabularx}{\textwidth}{L{2.7cm} L{1.4cm} L{1.65cm} L{1.75cm} c c c}
        \toprule
        \textbf{Dataset} & \textbf{Architecture} & \textbf{Pre-processing} &
        \textbf{Post-processing} & \textbf{IoU} & \textbf{Dice} &
        \textbf{95\% CI} \\
        \midrule
        \texttt{acoustic\_scattering} & [\,] & [\,] & [\,] & [\,] & [\,] & [\,] \\
        \texttt{active\_matter} & [\,] & [\,] & [\,] & [\,] & [\,] & [\,] \\
        \texttt{euler\_multi\_quadrants} & [\,] & [\,] & [\,] & [\,] & [\,] & [\,] \\
        \texttt{gray\_scott\_reaction\_diffusion} & [\,] & [\,] & [\,] & [\,] & [\,] & [\,] \\
        \texttt{helmholtz\_staircase} & [\,] & [\,] & [\,] & [\,] & [\,] & [\,] \\
        \texttt{planetswe} & [\,] & [\,] & [\,] & [\,] & [\,] & [\,] \\
        \texttt{rayleigh\_benard} & [\,] & [\,] & [\,] & [\,] & [\,] & [\,] \\
        \texttt{shear\_flow} & [\,] & [\,] & [\,] & [\,] & [\,] & [\,] \\
        \texttt{turbulent\_radiative\_layer\_2D} & [\,] & [\,] & [\,] & [\,] & [\,] & [\,] \\
        \texttt{viscoelastic\_instability\_v2} & [\,] & [\,] & [\,] & [\,] & [\,] & [\,] \\
        \bottomrule
    \end{tabularx}
\end{table*}

Table~\ref{tab:cross_dataset} provides a cross-dataset summary. Because the
physical targets may differ between datasets, values should be compared
directly only when the label semantics and evaluation protocol are
sufficiently equivalent. Otherwise, report per-dataset ranks or normalized
improvements relative to each dataset's baseline.

\begin{table*}[t]
    \centering
    \caption{Cross-dataset summary of the selected segmentation approach.}
    \label{tab:cross_dataset}
    \small
    \begin{tabularx}{\textwidth}{L{3.1cm} L{2.1cm} L{2.1cm} c c Y}
        \toprule
        \textbf{Dataset} & \textbf{Selected architecture} &
        \textbf{Selected P / Q} & \textbf{Test Dice} &
        \textbf{Relative gain (\%)} & \textbf{Generalization notes} \\
        \midrule
        \texttt{acoustic\_scattering} & [\,] & [\,] & [\,] & [\,] & [\,] \\
        \texttt{active\_matter} & [\,] & [\,] & [\,] & [\,] & [\,] \\
        \texttt{euler\_multi\_quadrants} & [\,] & [\,] & [\,] & [\,] & [\,] \\
        \texttt{gray\_scott\_reaction\_diffusion} & [\,] & [\,] & [\,] & [\,] & [\,] \\
        \texttt{helmholtz\_staircase} & [\,] & [\,] & [\,] & [\,] & [\,] \\
        \texttt{planetswe} & [\,] & [\,] & [\,] & [\,] & [\,] \\
        \texttt{rayleigh\_benard} & [\,] & [\,] & [\,] & [\,] & [\,] \\
        \texttt{shear\_flow} & [\,] & [\,] & [\,] & [\,] & [\,] \\
        \texttt{turbulent\_radiative\_layer\_2D} & [\,] & [\,] & [\,] & [\,] & [\,] \\
        \texttt{viscoelastic\_instability\_v2} & [\,] & [\,] & [\,] & [\,] & [\,] \\
        \bottomrule
    \end{tabularx}
\end{table*}

\subsection{Statistical reporting}
\label{subsec:statistics}

For each reported configuration, provide the mean and standard deviation
across independent random seeds and/or test trajectories, and confidence
intervals where appropriate. State the statistical unit used for uncertainty
estimation. Avoid treating temporally adjacent frames from the same trajectory
as independent observations. Report class-wise scores for imbalanced
segmentation tasks, together with macro-averaged scores. The test set must not
be used to select thresholds, post-processing parameters or model
hyperparameters.


\section{Conclusions}\label{sec:conclusions}

[Synthesis of the main takeaways of the review, highlighting the strengths and limitations of the surveyed methods.]

[Paragraph outlining key open challenges and concrete directions for future research (e.g., real-time feasibility, sparse sensor networks, concurrent events, and advanced learning paradigms).]

%% ============================================================================
%% SECOES FINAIS OBRIGATORIAS (BACKMATTER — ORDEM IDENTICA AO ARTIGO)
%% ============================================================================
\section*{CRediT authorship contribution statement}
\textbf{[First Author]:} Writing -- review \& editing, Writing -- original draft, Supervision, Resources, Project administration, Investigation, Funding acquisition.
\textbf{[Second Author]:} Writing -- original draft, Validation, Methodology, Investigation, Formal analysis, Data curation, Conceptualization.
\textbf{[Third Author]:} Writing -- review \& editing, Supervision, Funding acquisition.

\section*{Declaration of competing interest}
The authors declare that they have no known competing financial interests or personal relationships that could have appeared to influence the work reported in this paper.

\section*{Data availability}
No data was used for the research described in the article.

%% Opcional (descomente se aplicavel):
%% \section*{Acknowledgements}
%% [Funding sources, grants, facilities, and collaborators.]

%% ============================================================================
%% REFERENCIAS (Estilo numerico Elsevier: elsarticle-num)
%% ============================================================================
%% OPCAO 1 (Recomendada com arquivo .bib):
%% \bibliographystyle{elsarticle-num}
%% \bibliography{references}

%% OPCAO 2 (Bibliografia manual embutida para compilacao imediata):
\begin{thebibliography}{99}

\bibitem{ref1}
A.~Author, B.~Author, C.~Author,
Title of the referenced journal article,
\textit{Results in Engineering} 26 (2025) 105568.
\url{https://doi.org/10.1016/j.rineng.2025.105568}

\bibitem{ref2}
D.~Author, E.~Author,
Title of another referenced article,
\textit{Mechanical Systems and Signal Processing} 220 (2024) 111608.

\bibitem{ref3}
F.~Author,
Title of a review or conference paper,
\textit{Computers \& Structures} 89 (23--24) (2011) 2093--2112.

\end{thebibliography}

\end{document}